\section{Progettazione}

\subsection{Use Case Template}
Di seguito vengono descritti i template per i casi d'uso indicati nello Use Case Diagram.
\subsubsection{Login}
\begin{table}[h!]
    \begin{tabularx}{\textwidth}{|c|X|}
    \hline
        \textbf{Use Case Template} \#1 & Login  \\ \hline
        \textbf{Descrizione} & L'utente accede al sistema attraverso le proprie credenziali \\ \hline
        \textbf{Livello} & User Goal \\ \hline
        \textbf{Attori} & Library User, Librarian o Library Administrator \\ \hline 
        \textbf{Pre Condizioni} & L'utente è nella pagina appropriata e non ha già fatto il login \\ \hline
        \textbf{Basic Flow} & \begin{enumerate}[nosep, leftmargin=* ,before=\vspace{-0.5\baselineskip}]
            \item L'utente inserisce il nome utente e password 
            \item L'utente invia le credenziali 
            \item Il sistema controlla le credenziali 
            \item Il sistema autentica l'utente, istanziandone il proprio controller 
        \end{enumerate} \\ \hline
        \textbf{Alternative Flow} & 3a. Se le credenziali sono errate il sistema restituisce un errore permettendo di ritentare  \\ \hline
        \textbf{Post Condizioni} & L'utente è autenticato. Viene istanziato il controller appropriato\\ \hline
        \textbf{Test} & 1,2,3 \\ \hline
    \end{tabularx}
    \caption{Login UCT}
    \label{tab:Login UCT}
\end{table}

\subsubsection{Registrazione}
\begin{table}[h!]
    \begin{tabularx}{\textwidth}{|c|X|}
    \hline
        \textbf{Use Case Template} \#1 & Registrazione  \\ \hline
        \textbf{Descrizione} & L'utente sprovvisto di credenziali si registra \\ \hline
        \textbf{Livello} & User Goal \\ \hline
        \textbf{Attori} & Library User, Librarian o Library Administrator \\ \hline 
        \textbf{Pre Condizioni} & L'utente è nella pagina appropriata e non ha fatto il login \\ \hline
        \textbf{Basic Flow} & \begin{enumerate}[nosep, leftmargin=* ,before=\vspace{-0.5\baselineskip}]
            \item L'utente inserisce nome, cognome, email, password 
            \item L'utente invia le credenziali 
            \item Il sistema controlla che non esista già un utente con l'email inserita 
            \item Il sistema crea l'account 
        \end{enumerate} \\ \hline
        \textbf{Alternative Flow} & 3a. Se l'utente esiste già restituisce un errore permettendo di ritentare  \\ \hline
        \textbf{Post Condizioni} & Viene creato l'account, permettendo all'utente di fare il login\\ \hline
        \textbf{Test} & 1,2,3 \\ \hline
    \end{tabularx}
    \caption{Registrazione UCT}
    \label{tab:Registrazione UCT}
\end{table}

\subsubsection{Reset password}
\begin{table}[H]
    \begin{tabularx}{\textwidth}{|c|X|}
    \hline
        \textbf{Use Case Template} \#1 & Reset password  \\ \hline
        \textbf{Descrizione} & Permette all'utente che si è dimenticato la propria password di recuperare il profilo reimpostando la password \\ \hline
        \textbf{Livello} & User Goal \\ \hline
        \textbf{Attori} & Library User, Librarian o Library Administrator \\ \hline 
        \textbf{Pre Condizioni} & L'utente è nella pagina appropriata, dopo aver ricevuto la possibiltà di recuperare la password confermando la propria identità \\ \hline
        \textbf{Basic Flow} & \begin{enumerate}[nosep, leftmargin=* ,before=\vspace{-0.5\baselineskip}]
            \item L'utente inserisce la nuova password, confermandola 
            \item L'utente invia le credenziali 
            \item Il sistema cambia la password con quella indicata 
        \end{enumerate} \\ \hline
        \textbf{Alternative Flow} & -  \\ \hline
        \textbf{Post Condizioni} & Viene resettata la password\\ \hline
        \textbf{Test} & 1,2,3 \\ \hline
    \end{tabularx}
    \caption{Reset password UCT}
    \label{tab:Reset password UCT}
\end{table}